\section{Supplementary support reductions and continuations}\label{app:supplementary}

The first subsection gives an elementary counterpart to the Aebi-based
exclusion. The later subsections distinguish necessary incidence data and
computational organization from the unresolved Diophantine questions.


\subsection{The elementary endpoint route}

For comparison, we retain the local endpoint argument and valuation
restrictions. In this subsection the configuration is primitive
and satisfies the full-magic interaction.

Throughout, $\supp(2d)=\{2,3,p\}$ with $p\neq3$ prime, and
$d/2=XYZ$ with $X=2^{\alpha}$, $Y=3^{\beta}$, $Z=p^{\gamma}$. The two
technical stages are stated first; the headline
Theorem~\ref{main:threeprime} subsumes both by a different global
argument, and its proof is independent of theirs.

\begin{proposition}[All-endpoint case, $p\equiv3\bmod8$]\label{prop:threeblock}
There is no primitive positive full-magic transversal configuration with
$\supp(2d)=\{2,3,p\}$ and $p\equiv3\pmod8$.
\end{proposition}

\begin{proof}
By the endpoint table of Section~\ref{sec:supportlaw}, $p\equiv3\bmod8$
forces all three indices to be endpoints at $p$, at $3$ (same residue
argument), and --- by the opposite-parity mechanism at $2$ --- at the
$2$-power as well. Each progression's equal-area Pythagorean triangle is
then primitive, and its legs are a unitary partition of the three whole
blocks $X,Y,Z$. Apart from the impossible split $\{1,XYZ\}$, only
$\{X,YZ\}$, $\{Y,XZ\}$, $\{Z,XY\}$ remain, and three pairwise-distinct
progressions must use all three. Factoring the hypotenuse equations of the
last two gives
\[
Y^2-1=2XZ,\qquad Z^2-1=2XY ,
\]
whose difference factors as $(Y-Z)(Y+Z+2X)=0$. Positivity forces $Y=Z$,
after which $Y(Y-2X)=1$, impossible for positive prime-power blocks.
\end{proof}

\begin{proposition}[One-interior law, $p\equiv5,7\bmod8$]\label{prop:oneinterior}
Let $\supp(2d)=\{2,3,p\}$ with $p\equiv5$ or $7\pmod8$. The all-endpoint
branch is impossible by the three-block argument, so $\gamma=1$ is
excluded, and for $\gamma\ge2$ exactly one index is interior at $p$. If
$p\equiv5\pmod8$ the centre progression is the interior one and exactly one
of
\[
v_p(X^2+Y^2)=2h,\qquad v_p(1+Y^2)=2h,\qquad v_p(1+X^2)=2h
\]
holds, with $h=v_p(s_0)\ge1$; if $p\equiv7\pmod8$ one outer progression is
interior and exactly one of the six ordered differences
\[
v_p(2X^2-Y^2),\; v_p(2Y^2-X^2),\; v_p(2-Y^2),\;
v_p(2-X^2),\; v_p(2Y^2-1),\; v_p(2X^2-1)
\]
equals $2h$, with $h=\min(a,\gamma-a)$. These are necessary branch laws,
not contradictions; their exact certificate covers all nine identities.
\end{proposition}

The nine residual branches of Proposition~\ref{prop:oneinterior} need not
be analyzed one by one: the headline theorem closes every case, including
$p\equiv1\bmod8$, at once.

\begin{proof}[An elementary proof for primitive exact support $\{2,3,p\}$]
By the endpoint table, every index is an endpoint at $3$ and at least two
indices are endpoints at $p$. Those two indices give two \emph{distinct
primitive} Pythagorean triangles of the same area $d/4$, whose area has
exactly three distinct prime factors. Writing $X=2^{\alpha}=2x$, the three
endpoint types satisfy
\[
YZ=x^2-1,\qquad Y^2-1=4xZ,\qquad Z^2-1=4xY .
\]
Combining the first with each of the others reduces both mixed pairs to
one equation
\[
W(W^2-1)=4x(x^2-1),
\]
where $W$ is the remaining odd block. The $2$-adic valuation of the right
side forces $W\equiv\pm1\pmod{2x}$, while monotonicity of $W(W^2-1)$ gives
$x<W<2x$; the only residue in that window is $W=2x-1$, which forces
$x=2$, $W=3$, and then the third block equals $1$ --- not a prime power in
the support. The remaining pair of endpoint types forces an equality of
powers of two distinct primes, equally impossible. Hence no primitive
positive full-magic transversal configuration has $\supp(2d)=\{2,3,p\}$ for any prime
$p\neq3$.
\end{proof}

\begin{remark}[Attribution and scope]
The headline theorem is the squareclass consequence of Aebi's complete
triangle classification proved in Section~\ref{sec:threeprime}. The endpoint
argument gives another proof of its primitive full-magic consequence,
using only the stated endpoint restrictions. The uniqueness assertion fails already at four
primes, as the difference-$840$ example shows. Section~\ref{sec:smoothsupport}
nevertheless closes particular four-prime supports by a separate complete
catalogue; arbitrary four-prime supports and global support-size control
remain open.
\end{remark}


\subsection{The difference-side taxonomy}

Consider a root-primitive positive full-magic transversal configuration.
For each progression put
\[
U_i=(t_i-r_i)/2,\qquad V_i=(t_i+r_i)/2,
\]
so $U_iV_i=d/2$ and $U_i^2+V_i^2=s_i^2$. Call its AP triangle
\emph{primitive} when $\gcd(U_i,V_i)=1$, and let $I$ be the set of vertices
whose triangles are not primitive. A prime is interior at $i$ precisely
when it divides both $U_i$ and $V_i$. By the support law in
Section~\ref{sec:supportlaw}, neither $2$ nor $3$ is interior, and each
other support prime is interior at at most one vertex. Thus
\[
\#\{\text{primitive AP triangles}\}
 =3-|I|\ge\max\{0,5-\omega(2d)\}.
\]
Indeed, each occupied vertex requires a distinct prime outside $\{2,3\}$.
The eight subsets of $\{0,1,2\}$ give eight labelled incidence cases,
without asserting that any case is realized by a configuration.

The allowed interior vertices for an odd prime are $\{0,1,2\}$,
$\varnothing$, $\{0\}$ and $\{1,2\}$ for residues $1,3,5,7\bmod8$,
respectively. An interior prime also requires $v_p(d)\ge2$. For a fixed
$d$, let $n_j$ count the primes $p\notin\{2,3\}$ with this exponent
condition and $p\equiv j\bmod8$. In the relaxation that retains only
these residue and exponent conditions, assigning distinct primes to all
three vertices is possible exactly when
\[
n_1+n_5\ge1,\qquad n_1+n_7\ge2,\qquad n_1+n_5+n_7\ge3.
\]
To see this, apply Hall's matching criterion to the vertices and eligible
primes. The displayed inequalities are its conditions for the central
singleton, the pair of outer vertices, and all three vertices. The two
outer vertices have identical neighbor sets, and every remaining subset
condition follows from those three. This is an equivalence for the
matching relaxation only; the Pythagorean and full-magic interaction
equations still have to hold.

\subsection{The outer-primitive reduction}

A valuation record organizes possible branches but does not produce the
required triangles. A further elementary reduction applies when the two
outer triangles are both primitive. Set
\[
G=\gcd(U_1,U_2),\qquad R=U_1/G,\qquad S=U_2/G.
\]
Since $\gcd(R,S)=1$ and $U_1V_1=U_2V_2$, there is a positive integer $B$
such that
\[
(U_1,V_1)=(GR,BS),\qquad (U_2,V_2)=(GS,BR).
\]
The two primitive gcd conditions make $G,B,R,S$ pairwise coprime.
Exactly one factor is even, and it is divisible by four because the even
leg of a primitive Pythagorean triangle is divisible by four.

Put $A=G^2+B^2$ and $C=R^2+S^2$. Exactly one of these norms is twice an
odd integer; divide that norm by two to obtain odd integers $A_0,C_0$.
The full-magic interaction gives
\[
2s_0^2=s_1^2+s_2^2
 =(G^2+B^2)(R^2+S^2),\qquad A_0C_0=s_0^2.
\]
Hence $A_0$ and $C_0$ have the same unique positive squarefree kernel.
Moreover, if an odd prime $q$ divides both, pairwise coprimality permits
division by $B,S$ modulo $q$, giving
\[
(G/B)^2\equiv(R/S)^2\equiv-1,\qquad
(s_1/(BS))^2\equiv2\pmod q.
\]
The last congruence uses $s_1^2=G^2R^2+B^2S^2$. Thus both $-1$ and $2$
are quadratic residues modulo $q$, and $q\equiv1\pmod8$.

The common-kernel conclusion uses the central interaction as well as the
two primitive outer triangles. Conversely, a common kernel does not
supply the two outer-square equations. These necessary reductions leave
the nondegenerate rational and integral point problems open.

\subsection{Where the open objects live}

The companion repository records three continuation boundaries in
\texttt{dossiers/README.md}. These are status summaries of further
arithmetic work. The shipped \texttt{scripts/verify.py} checks exact
main-paper identities and the integrity and content of specified frozen
computer-algebra outputs; it does not reproduce the deeper models and
computations summarized in these dossiers. The mathematical conclusions
depend on the proofs and the completeness statements of the imported
computations, not on successful output checks alone.

\begin{itemize}
\item \textbf{Residual cover}
(\texttt{dossiers/}\allowbreak\texttt{residual\_cover/README.md}).
The continuation from the area-$30/60$ covering classes of
Section~\ref{sec:frontier} records a residual cover $q_R$ over
$K=\Q(\sqrt{218})$. Neither a $K$-point nor emptiness is established.
The recorded alternatives for the associated elliptic curve are rank two
if a point exists, and rank one
with a nontrivial divisible Tate--Shafarevich class if the cover is empty.

\item \textbf{Odd lift}
(\texttt{dossiers/}\allowbreak\texttt{odd\_lift\_n5/README.md}).
The coefficient-$3$ result excludes nondegenerate rational
specializations of the known generic coefficient-$3$ family satisfying
both base reconstructions; it does not classify specialized
Mordell--Weil points outside the generic subgroup.
For coefficient $5$, the reduction has two genus-two covers over
$\Q(\sqrt2)$. One is empty on the required real component; the surviving
cover has a complete fake two-Selmer set of four classes, none eliminated.
The associated degree-$30$ pair-field computation gives a rank-two
lattice saturated at $2,3,5,7$, without a proved rank upper bound or finite
index in the full Mordell--Weil group. A complete list of points with the
required rational coordinate and a uniform result for larger odd
coefficients remain open.

\item \textbf{Degree-$90$ $S$-class}
(\texttt{dossiers/}\allowbreak\texttt{degree90\_sclass/README.md}).
The marked ideal class $[A_R]$ in a degree-$90$ $S$-class group remains
undecided. The dossier records a field construction and structural
reductions, but bounded short-vector and elementary-unit searches
produced no witness. Their failure does not prove that the class is
nontrivial; the remaining decision requires a global witness or a
rigorous obstruction.
\end{itemize}

The limits of earlier approaches are recorded in
\texttt{dossiers/}\allowbreak\texttt{EXHAUSTED\_ROUTES.md}.
These continuations do not decide every branch as the prime support
varies. For a fixed support, the two-term $S$-unit route of
Section~\ref{sec:smoothsupport} gives a finite reduction; the uniform
all-support problem remains open.
